\documentclass[10pt,a4paper]{amsart}
\usepackage[margin=19mm]{geometry}
\usepackage{amsmath,amssymb,mathtools}
\usepackage[scr=rsfs,cal=euler]{mathalfa}
\usepackage{xfrac}
\usepackage[unicode,hidelinks,bookmarks=false]{hyperref}
\newcommand{\ie}{i.e.\ }
\newcommand{\cf}{cf.\ }
\newcommand{\resp}{respectively\ }
\newcommand{\Subsection}{\subsection}
\providecommand{\nobreakdash}{\nobreak\hbox{-}}
\setlength{\emergencystretch}{2em}
\multlinegap0pt
\usepackage{tikz}
\usepackage[shortlabels]{enumitem}
\usepackage{amssymb}
\newtheorem{lemma}{Lemma}
\newtheorem{theorem}{Theorem}
\newcommand{\ZZ}{\mathbb{Z}}
\title{Integral bases for the second degree cohomology of 4-dimensional toric orbifolds}
\author{Tseleung So, Jongbaek Song}
\date{Source: arXiv:2604.02562v1 (CC BY 4.0)}
\begin{document}
\maketitle

\noindent\emph{Source and licence.} Integral bases for the second degree cohomology of 4-dimensional toric orbifolds. Version arXiv:2604.02562v1 (CC BY 4.0), distributed by arXiv under the Creative Commons Attribution 4.0 International licence, http://creativecommons.org/licenses/by/4.0/, as recorded in the arXiv licence metadata for this exact version and shown on the abstract page https://arxiv.org/abs/2604.02562v1. Full licence notice: \texttt{licenses/arxiv-2604.02562-CC-BY-4.0.txt}.\par\smallskip

\noindent\textit{Editorial scope.} Counted unit: the article's theorem thm\_main\_lattice\_intersection (source0.tex lines 467--475). The article's complete original proof of it is delivered with it (source0.tex lines 477--550, one \texttt{proof} environment of 74 lines). No mathematics is written, repaired or re-derived: the statement and the proof are the article's own lines, in the article's own notation, and no retained line is substituted or re-flowed.  Retained uncounted as the source-local context the proof needs: lines 446--451.  Nothing was omitted from any retained span, so the package carries no omission record.  No retained line contains a citation, so the wrapper carries no bibliography and no bibliography entry.  No retained line contains a figure environment, an image-inclusion command or drawing code, so no image is delivered and the package declares no asset dependency.  Editorial formatting only, in addition: an AMS \texttt{amsart} preamble that restates the article's own declarations and macros used by the retained lines, namely the environments \texttt{lemma}, \texttt{theorem} on separate counters, together with the standard packages those lines need.  The retained environments print in this document as Lemma 1, Theorem 1 in order of appearance, whereas the article prints the same items with the numbering of its own full document.  The complete native source of the article as distributed in the arXiv e-print for this exact version is bundled unchanged as \texttt{sources/197745/source0.tex}.
\begin{lemma}\label{lem_w_m_0}
\begin{enumerate}
    \item If $(t_1, \dots, t_m)\in K$, then $w_m=0$. 
    \item The restriction $\phi|_K \colon K \to \ZZ^m$ is injective. 
\end{enumerate}
\end{lemma}

\begin{theorem}\label{thm_main_lattice_intersection}
For the lattices $L_i$ defined above for $i=1, \dots, m$,  it follows that 
\[
L_1\cap \cdots \cap L_m = \mathbb{Z}\left<\mathbf{a}\right> \oplus \mathbb{Z}\left<\mathbf{b}\right>\oplus \phi(K).
\]
In particular, if $\{\kappa_1, \dots, \kappa_{m-2}\}$ is a $\mathbb{Z}$-basis of $K$, then
$\{ \mathbf{a}, \mathbf{b}, \phi(\kappa_1), \dots,  \phi(\kappa_{m-2})\}$
is a $\mathbb{Z}$-basis of $L_1\cap\cdots\cap L_m$.
\end{theorem}

\begin{proof}
We first claim that 
\begin{equation}\label{eq_int_equal_sum}
L_1\cap \cdots \cap L_m=\mathbb{Z}\left<\mathbf{a}\right> + \mathbb{Z}\left<\mathbf{b}\right>+ \phi(K).
\end{equation}
We start with the inclusion  $L_1\cap \cdots \cap L_m\subseteq\mathbb{Z}\left<\mathbf{a}\right> + \mathbb{Z}\left<\mathbf{b}\right>+ \phi(K)$.
Note that $\mathbf{x}=(x_1, \dots, x_m)\in L_i$ if and only if 
\begin{equation}\label{eq_coord_condition}
x_i=a_ip_i+b_iq_i,\quad \text{and} \quad x_{i+1}=a_{i+1}p_i+b_{i+1}q_i
\end{equation}
for some $p_i, q_i\in \ZZ$. Therefore, for $\mathbf{x}=(x_1, \dots, x_m)\in L_1\cap\cdots \cap L_m$, there exist
$p_1, q_1, p_2, q_2, \dots, p_m, q_m\in \ZZ$ 
such that \eqref{eq_coord_condition} holds for each $i=1, \dots, m$. Here, indices are taken modulo $m$, that is, $(a_{m+1}, b_{m+1})=(a_1, b_1)$. Hence, we get 
\[
p_ia_i+q_ib_i= p_{i-1}a_i + q_{i-1}b_i,
\]
for each $i=1, \dots, m$, where we set $(p_0, q_0)=(p_m, q_m)$. 

Since $\lambda_i=(a_i, b_i)$ is primitive for each $i=1,\ldots,m$, there exists $t_i\in\ZZ$ such that
\begin{equation}\label{eq_pq_ab}
(p_i-p_{i-1}, q_i- q_{i-1}) = t_i (-b_i, a_i)
\end{equation}
Summing Equation~\eqref{eq_pq_ab} over $i=1, \dots, m$, we get 
\[
0=\sum_{i=1}^m (p_i-p_{i-1})  = - \sum_{i=1}^m b_it_i \quad \text{and} \quad 
0=\sum_{i=1}^m (q_i-q_{i-1})  =  \sum_{i=1}^m a_it_i,
\]
implying that $(t_1, \dots, t_m)\in K$. 
Similarly, summing the first $i$ equations in~\eqref{eq_pq_ab} gives
\[
p_i=p_m - \sum_{j=1}^i b_jt_j \quad \text{and} \quad q_i=q_m + \sum_{j=1}^i a_jt_j.
\]
Therefore, we have for each $i=1, \dots, m$,
\begin{align*}
x_i&=p_ia_i+q_ib_i\\
&=\left(p_m - \sum_{j=1}^i b_jt_j\right) a_i + \left( q_m + \sum_{j=1}^i a_jt_j \right) b_i \\
&=p_ma_i+q_mb_i + \sum_{j=1}^i (a_jb_i - a_ib_j) t_j\\
&= p_ma_i+q_mb_i + \sum_{j=1}^{i-1} (a_jb_i - a_ib_j) t_j.
\end{align*}
This establishes that $\mathbf{x}=(x_1, \dots, x_m)\in \mathbb{Z}\left<\mathbf{a}\right> + \mathbb{Z}\left<\mathbf{b}\right>+ \phi(K)$. 

We now prove the reverse inclusion $\mathbb{Z}\left<\mathbf{a}\right> + \mathbb{Z}\left<\mathbf{b}\right>+ \phi(K)\subseteq L_1\cap \cdots \cap L_m$. For an element $\alpha \mathbf{a} + \beta \mathbf{b} + \phi(t)$ with $\alpha, \beta\in \mathbb{Z}$ and $t=(t_1, \dots, t_m) \in K$, we set 
\[
p_i=\alpha - \sum_{j=1}^i b_jt_j \quad \text{and} \quad q_i=\beta + \sum_{j=1}^i a_jt_j.
\]
Observe that $p_m= \alpha$ and $q_m= \beta$ as $t\in K$.
Then, writing $x_i$ as the $i$-th coordinate of $\alpha \mathbf{a} + \beta \mathbf{b} + \phi(t)$, we have 
\[
x_i = \alpha a_i + \beta b_i + \sum_{j=1}^{i-1}(a_jb_i-a_ib_j)t_j = \alpha a_i + \beta b_i + \sum_{j=1}^{i}(a_jb_i-a_ib_j)t_j  =  a_ip_i+b_iq_i
\]
for $i=1,\dots, m-1$. Furthermore, it follows that 
\begin{equation*}
% \label{eq_m_coor}
\begin{split}
x_m& =\alpha a_m + \beta b_m + \sum_{j=1}^{m-1}(a_jb_m-a_mb_j)t_j\\
&= \alpha a_m + \beta b_m + \sum_{j=1}^{m-1} \det (\lambda_j, \lambda_m)t_j  \\
%&= \alpha a_m + \beta b_m + \det ( t_1\lambda_1 + \cdots + t_{m-1}\lambda_{m-1}, \lambda_m)\\
&= \alpha a_m + \beta b_m 
\end{split}
\end{equation*}
where the last equality follows from Lemma \ref{lem_w_m_0}-(1). A similar computation gives the second equation in \eqref{eq_coord_condition}. Therefore
$\alpha \mathbf{a} + \beta \mathbf{b} + \phi(t)\in L_1 \cap \cdots \cap L_m$.
% and hence Equation~\eqref{eq_int_equal_sum} holds.

It remains to show that the right hand side of \eqref{eq_int_equal_sum} is indeed a direct sum $\ZZ\langle\textbf{a}\rangle\oplus\ZZ\langle\textbf{b}\rangle\oplus\phi(K)$. 
Assume that $\alpha \mathbf{a} + \beta \mathbf{b} + \phi(t)=0$ for some $\alpha, \beta\in \ZZ$ and $t\in K$. Note that if $\phi(t)=(w_1, \dots, w_m)$, then $w_1=0$ by convention and $w_m=0$ by Lemma \ref{lem_w_m_0}-(1). Then, we have 
\[
\alpha a_1+ \beta b_1 = \alpha a_m + \beta b_m = 0
\]
which implies that $\alpha=\beta=0$ because $\lambda_1=(a_1, b_1)$ and $\lambda_m=(a_m, b_m)$ are linearly independent. Hence, the result follows. 


Finally, the second claim follows directly from Lemma \ref{lem_w_m_0}-(2). This completes the proof.
\end{proof}

\end{document}
